% ---------------------------------------------------------------------------
% MASTER RESUME - NEVER SUBMIT THIS FILE.
%
% This is the printable superset: every role, every project, every bullet
% variant, all in one place. Run it to 2-3 pages and let it grow.
%
% Bracketed tags mark which base each bullet is written for:
%   [emb]  embedded / controls
%   [sim]  simulation / robotics / test
%   [swe]  general software / backend
%   [lead] leadership / ownership framing
%
% master/bullet_bank.md is the source of truth for wording and confidence.
% When you change a bullet there, mirror it here.
% ---------------------------------------------------------------------------
\documentclass[10pt,letterpaper]{article}
\usepackage{resume}

\pagestyle{plain}   % master is multi-page; keep page numbers

\newcommand{\vtag}[1]{\textbf{[#1]}~}

\begin{document}
\bodyfont

\resumeheader

\begin{center}
{\fontsize{9}{10}\selectfont\textbf{MASTER COPY - NOT FOR SUBMISSION}}
\end{center}
\vspace{4pt}

\educationsection

\sectionheader{SKILLS INVENTORY (full, unordered)}

\techline{\textbf{Languages}: C, C++, Python, TypeScript/JavaScript, SQL, Java, Assembly}
\techline{\textbf{Controls}: Fixed-point PI/PID, anti-windup, saturation, slew-rate limiting, control-loop tuning, state estimation, sensor fusion, Bode/Nyquist}
\techline{\textbf{Embedded}: TTC 60 VCU, STM32H7, PX4 flight stack, real-time firmware, CAN, SPI, microcontroller I/O}
\techline{\textbf{Simulation \& Test}: Software-in-the-loop, closed-loop simulation, scripted replay, regression suites, parameter sweeps, pytest}
\techline{\textbf{Robotics \& Perception}: Visual-inertial odometry, IMU fusion, OpenCV, real-time vision pipelines}
\techline{\textbf{Web \& APIs}: Flask, FastAPI, React, Node.js/Express, REST/JSON, input validation, sessions/JWT basics}
\techline{\textbf{Data}: NumPy, pandas, Matplotlib, SQLite, SQL Server, schema design, CSV/JSON telemetry pipelines}
\techline{\textbf{Quality \& DevOps}: pytest, Jest, GitHub Actions, Docker, structured logging, GDB, Git/GitHub, Postman}
\techline{\textbf{Domain Tools}: TASKING, TTC Downloader, BusMaster, MoTeC i2 Pro, Pygame}

\sectionheader{EXPERIENCE - ALL BULLETS}

\roleline{Software Lead}{Jun 2026 - Present}
\orgline{\spartanracing}
\begin{itemize}
  % Rewritten 2026-08-16 from the Aug 12 2026 Electronics \& Software meeting deck.
  % Present tense throughout: most of this is in-progress, not shipped. See bullet_bank.md.
  \item \vtag{lead} Lead the software organization for a Formula SAE electric racecar, directing \textbf{12 engineers} across \textbf{10+ concurrent projects} spanning a custom vehicle control unit, torque vectoring, traction control, derating, launch control, plant modeling, BMS firmware, and telemetry
  \item \vtag{emb} Leading an in-house vehicle control unit program replacing a vendor VCU with a custom STM32H755 controller, with a 10 ms deterministic control task and a watchdog that reboots the controller if that task hangs
  \item \vtag{emb} Directing a seam-based port of a \textbf{5,400-line} production VCU codebase, fixing the sensor and CAN layer as the interface so everything below it could be rewritten for new silicon while control logic above carried over unmodified --- \textbf{4,100 lines} migrated to date
  \item \vtag{emb} Built the CAN driver for the custom VCU, sending and receiving on both controller buses, and specified the analog front-end changes needed for bring-up on the car
  \item \vtag{emb} Started a four-motor torque vectoring program, defining a phased plan from an R\&D minimum viable controller through motor and track simulation, with traction control layered on top of it
  \item \vtag{emb} Scoped a traction control loop for a four-motor car --- per-wheel slip from wheel speed against an IMU-integrated speed estimate, converted to per-wheel torque limits handed to torque vectoring and clamped to the 80 kW rule on a \textbf{5 ms} cycle
  \item \vtag{sim} Leading a model-, software-, and hardware-in-the-loop effort building battery, BMS, motor, and inverter plant models so VCU control code can be validated against a simulated vehicle before it reaches hardware
  \item \vtag{sim} Set code standards and an integration plan across a \textbf{7-engineer} modeling team so independently developed subsystem models compose into a single testbench
  \item \vtag{sim} Built a competition points model relating average lap time and efficiency score to points delta against the benchmark team, identifying a \textbf{59.25 s} lap target at a \textbf{40.73}-point efficiency score, and derived MoTeC i2 Pro channels to automate the underlying power-limit and regen analysis
  \item \vtag{swe} Standing up a build and release pipeline for VCU firmware, specifying the compute, storage, and CI infrastructure it needs and driving sponsorship outreach to \textbf{120+} hardware, cloud, and developer-infrastructure companies to fund it
  \item \vtag{emb} Leading firmware for a custom battery management system, covering cell monitoring, high-voltage and current sense over an ADBMS2950 front end, and the state-of-charge estimation built on it
  \item \vtag{emb} Leading dashboard firmware for a new three-CAN STM32 driver display, carrying the existing telemetry, energy-budget, and alerting features onto the new hardware
  \item \vtag{emb} Overseeing a thermal derating model heading to a first release as the season's baseline configuration, including GPU-accelerated sweeps and validation of VCU derating code through software- and hardware-in-the-loop
  \item \vtag{emb} Driving wireless firmware flashing onto the race car for track testing, taking the CAN flashing work from bench tooling to a modem-backed link usable at the test site
  \item \vtag{swe} Run a weekly technical review cadence where every project reports written completed and to-do status against plan, and authored onboarding documentation and a software curriculum for new members
\end{itemize}

\entryspace

\roleline{Software Engineer Intern}{Feb 2026 - Present}
\orgline{Argus Defense \textbar{} San Jose, CA}
\begin{itemize}
  \item \vtag{emb} Tuned PID flight-control loops on a PX4 flight controller to reduce vertical and horizontal position drift in autonomous flight
  \item \vtag{emb} Integrated and tuned visual-inertial odometry for onboard state estimation, improving position-hold accuracy
  \item \vtag{emb} Developed and validated embedded flight-control firmware in C on PX4, iterating control gains against logged flight telemetry
  \item \vtag{sim} Tuned PX4 PID flight-control loops to reduce vertical and horizontal drift during autonomous position hold
  \item \vtag{sim} Integrated and tuned visual-inertial odometry for onboard state estimation and evaluated performance against logged telemetry
  \item \vtag{sim} Developed embedded flight-control changes in C and iterated gains through repeatable flight tests and log review
  \item \vtag{swe} Developed and validated flight-control software in C on PX4, using structured telemetry logs to reproduce issues and evaluate changes
  \item \vtag{swe} Integrated visual-inertial odometry and iterated configuration through repeatable tests and documented results
\end{itemize}

\entryspace

\roleline{Software Controls Engineer}{Jul 2025 - Jun 2026}
\orgline{\spartanracing}
\begin{itemize}
  \item \vtag{emb} Designed a lap-adaptive energy-management algorithm that sets a power setpoint from cumulative energy versus budget, with the controller holding actual power within \textbf{0.01\%} of setpoint
  \item \vtag{emb} Developed embedded power-limiting firmware with a PI controller on the VCU, maximizing motor output while enforcing FSAE's 80 kW limit
  \item \vtag{emb} Built CAN diagnostics firmware for live power and torque frames, status bits, and error counters for real-time BusMaster monitoring
  \item \vtag{sim} Implemented a torque-to-power controller with fixed-point PI control, saturation, anti-windup, and slew limits for deterministic VCU execution
  \item \vtag{sim} Built Python telemetry parsers and a software-in-the-loop replay harness to validate controller changes against recorded CSV logs
  \item \vtag{sim} Planned and executed on-vehicle tests, measured response and latency, and generated plots for parameter tuning and regression review
  \item \vtag{swe} Developed Python telemetry pipelines to parse vehicle logs, calculate controller response metrics, and generate plots for tuning and validation
\end{itemize}

\entryspace

\roleline{Software Intern}{Aug 2024 - Jun 2025}
\orgline{\spartanracing}
\begin{itemize}
  \item \vtag{emb} Built deterministic entry and exit state logic for power limiting from measured power, requested torque, and threshold flags
  \item \vtag{emb} Tuned the PI torque controller used for power limiting to within \textbf{2\%} error, smoothing torque transitions to preserve drivability
  \item \vtag{sim} Extracted a desktop simulator and regression suite to reproduce field scenarios, run controller sweeps, and prevent regressions
  \item \vtag{sim} Tuned PI torque control and deterministic state transitions using logged power, torque, and threshold data
  \item \textbf{Results:} 1st in endurance at MIS EV 2025
\end{itemize}

\entryspace

\roleline{Software Development Intern}{Sep 2024 - Oct 2024}
\orgline{Zip Intelligence \textbar{} Tracy, CA}
\begin{itemize}
  \item \vtag{swe} Implemented frontend validation and clearer error states, then paired with engineers to harden REST endpoints for status codes and edge cases
  \item \vtag{swe} Added SQL Server update scripts and structured logs to speed reproduction and issue triage
  \item \vtag{swe} Took ownership of a small feature from ticket through verified release, shipping via pull requests with setup and change documentation
\end{itemize}

\entryspace

\roleline{Data Analyst}{Jan 2024 - Aug 2024}
\orgline{Sailfish \textbar{} Tracy, CA}
\begin{itemize}
  \item \vtag{swe} Built structured dashboards, weekly summaries, and maintainable QA notes while streamlining data entry and reporting a 32\% improvement in tracked metrics
\end{itemize}

\sectionheader{PROJECTS - ALL BULLETS}

\roleline{Formula SAE EV Dashboard}{Feb 2026 - Present}
\orgline{\spartanracing}
\techline{C, STM32H7, CAN, SPI}
\begin{itemize}
  \item \vtag{emb} Built a real-time CAN telemetry display for HV voltage, throttle, cell temperature, and power budget on an LCD
  \item \vtag{emb} Designed an energy bar and LED alert system giving the driver lap-by-lap feedback on energy budget
\end{itemize}

\entryspace

\roleline{Formula SAE EV Vehicle Simulation}{Jul 2025}
\orgline{\spartanracing}
\techline{Python, Pygame, Matplotlib, vehicle dynamics, PID control, CSV logging}
\begin{itemize}
  \item \vtag{emb} Built a physics-based FSAE EV simulation to model acceleration using vehicle mass, drag, gearing, and motor torque
  \item \vtag{emb} Implemented PID-based power limiting to the FSAE 80 kW rule and logged runs to CSV for repeatable tuning and validation
  \item \vtag{sim} Built a physics-based vehicle simulator for torque, drag, acceleration, and steering to evaluate control changes before hardware tests
  \item \vtag{sim} Added headless batch runs, seeded scenarios, CSV summaries, and plots for repeatable offline comparisons
\end{itemize}

\entryspace

\roleline{SR-Wireless-CAN --- Wireless CAN Telemetry Platform}{2026}
\orgline{\spartanracing --- \textit{backend lead}}
\techline{Python, FastAPI, WebSockets, python-can, React, SQLite, Docker Compose}
\begin{itemize}
  \item \vtag{emb} Proposed and implemented wireless firmware flashing for the TTC 60 VCU, decoding the flashing sequence from CAN trace captures and reimplementing it in Python over the bus
  \item \vtag{swe} Led backend development for a wireless CAN telemetry platform, building a FastAPI and WebSocket service that streams live vehicle telemetry to a React dashboard with SQLite persistence
  \item \vtag{swe} Implemented structured DBC upload and merge with automatic conflict resolution, and a bounded search window that keeps large signal catalogs responsive
  \item \vtag{swe} Containerized the backend, frontend, and database with Docker Compose for one-command setup across team machines
  \item \vtag{sim} Added replay of recorded \texttt{.trc} CAN traces plus simulated and hardware (\texttt{can0}) capture modes, so runs could be reproduced and re-examined off the vehicle
  \item \vtag{sim} Built table, graph, and value/condition monitoring modes for arbitrary CAN signals, with live run observation shared across devices over WebSockets
  \item \vtag{emb} Developed wireless CAN capture and diagnostics tooling with DBC-driven signal decoding, replacing tethered bench monitoring during vehicle testing
\end{itemize}

\entryspace

\roleline{Posture Detection System}{2025}
\techline{Python, OpenCV, real-time vision, telemetry}
\begin{itemize}
  \item \vtag{sim} Developed a webcam pose-detection pipeline with overlays, smoothing, debounce logic, and configurable alert thresholds
  \item \vtag{sim} Tracked FPS and latency, recorded sessions, and exported labeled frames and CSV metrics for offline analysis
\end{itemize}

\entryspace

\roleline{Tickr --- AI-Coached Paper-Trading Game}{2025}
\techline{TypeScript/JavaScript, React, flat-file datastore, responsive UI}
\begin{itemize}
  \item \vtag{swe} Built registration, sign-in, game, and history flows for a scenario-driven trading game with an AI coach
  \item \vtag{swe} Implemented quest and prompt logic, client-side state management, and a mobile-responsive interface
\end{itemize}

\entryspace

\roleline{Scent Showdown}{2025}
\techline{Flask, HTML/CSS/JavaScript, SQLite, REST/JSON}
\begin{itemize}
  \item \vtag{swe} Built a matchup and leaderboard application with Flask routes, templates, a compact SQLite schema, input validation, and an admin view
  \item \vtag{swe} Added JSON endpoints, seed and export scripts, deployment setup, and a maintainer-focused README
\end{itemize}

\entryspace

\roleline{Educational Trading Bot}{---}
\techline{TO FILL IN - see master/project\_inventory.md}
\begin{itemize}
  \item Listed in the job-search plan but no bullets exist yet in any resume version. Write two bullets and add them here and to the bullet bank.
\end{itemize}

\end{document}
